The bidimensional steady compressible boundary layer equation for with zero
streamwise pressure gradient and non-body forces for a one step chemical
reaction can be writem for  conservation of mass as   

\begin{equation} 
\frac{\partial \rho^* u^* }{\partial x^*}+
\frac{\partial \rho^* v^* }{\partial y^*}+
=0
\end{equation}

Considering a single step chemical reaction,

\begin{equation}
\nu_F F + \nu_O O \rightarrow \nu_P P
\end{equation}

where $\nu_F$ and $\nu_O$
are the stiochiometric coeficients, for the 
fuel and oxidizer, respectively. 
Considering $\nu_F=1$ and $\nu_O=n$, indicating that 
$n$ moles of oxidizer are
required to burn one mole of fuel, producing $\nu_p=(1+n)$ moles of product,

\begin{equation}
F + n O \rightarrow (1+n)P
\end{equation}

where $n$ is the stiochiometric 
mass ration between the fuel and oxidizer, which can be expressed as

\begin{equation}
n=\frac{W_F}{W_O}\frac{\nu_O}{\nu_F}
.
\end{equation}

Where $W_F$ and $W_O$ are the molecular 
weights of the fuel and oxidizer, respectively.

Then the oxidizer mass consevation equation is 

\begin{equation} 
\rho^* u^*\frac{\partial Y_{O}^*}{\partial x^*}+
\rho^* v^*\frac{\partial Y_{O}^*}{\partial y^*}
=
\frac{\partial}{\partial y^*}
\left(\rho^* D_{O}^{*}\frac{\partial Y_O^{*}}{\partial y^{*}}\right)
-
\frac{\dot{\omega}}{n}
%\frac{\partial\rho D^m \tfrac{\partial\rho Y^m u_i}{\partial x_i}u_i}{\partial x_i}=
,
\end{equation} 

the fuel mass conservation

\begin{equation} 
\rho^* u^*\frac{\partial Y_{F}^*}{\partial x^*}+
\rho^* v^*\frac{\partial Y_{F}^*}{\partial y^*}
=
\frac{\partial}{\partial y^*}
\left(\rho^* D_{F}^{*}\frac{\partial Y_F^{*}}{\partial y^{*}}\right)
-
\dot{\omega}
%\frac{\partial\rho D^m \tfrac{\partial\rho Y^m u_i}{\partial x_i}u_i}{\partial x_i}=
\end{equation} 

and the product conservation 

\begin{equation} 
\rho^* u^*\frac{\partial Y_{P}^*}{\partial x^*}+
\rho^* v^*\frac{\partial Y_{P}^*}{\partial y^*}
=
\frac{\partial}{\partial y^*}
\left(\rho^* D_{P}^{*}\frac{\partial Y_P^{*}}{\partial y^{*}}\right)
+
(n+1)
\dot{\omega}
%\frac{\partial\rho D^m \tfrac{\partial\rho Y^m u_i}{\partial x_i}u_i}{\partial x_i}=
.
\end{equation} 

The momemtum streamwise equation 


\begin{equation} 
\rho u^*\frac{\partial u^*}{\partial x^*}+
\rho v^*\frac{\partial u^*}{\partial y^*}=-
\frac{\partial}{\partial y^*}
\left(\mu^*\frac{\partial u^*}{\partial y^*}\right)
.
\end{equation} 

The entalphy $h=e+p/\rho$ form of the energy equation:

\begin{equation}
%\frac{\partial}{\partial t}(\rho h)+
\rho^* u^* \frac{\partial h^*}{\partial x^*}+
\rho^* v^* \frac{\partial h^*}{\partial y^*}=
\mu^*\frac{\partial^2 u^*}{\partial (y^*)^2}+
\frac{\partial}{\partial y^*}
\left(
    k^*\frac{\partial T^*}{\partial y^*}
\right)
-
\frac{\partial}{\partial y^*}
\left(
        \sum_i
        h_{si}
        \rho 
        D_i
        \frac{\partial Y_i^*}{\partial y^*}
\right)
\end{equation} 


where $h_{si}$ is the sensible enthalpy 
for the specie $i$ in a multicomponent misture.  
For multiples species $h_s=\sum Y_ih_{si}(T)$
with  

\begin{equation} 
h_{si}^*
=
\int_{T_0^*}^{T^*}
c_{pi}^*dT^*
\end{equation} 

and the formation enthalpy of the misture  
can be defined 

\begin{equation} 
   h_{f}^\circ
   =
   \sum
   h_{fi}^{\circ*}
   .
\end{equation} 

Enthalpy can divide explicitly in two parts 

\begin{equation} 
    h^*=h_s^*+h_f^{\circ*}
\end{equation} 

for multiple species 

\begin{equation} 
    h^*=
    \sum_i
    \int_{T_0^*}^{T^*}
    c_{pi}^*dT^
    +
    \sum_i h_{fi}^{*\circ}
\end{equation} 


Taking the material derivative

\begin{equation} 
\begin{split}
    \frac{D h^*}{D t^*}
    =
    \frac{D h_s^*}{D t^*}
    +
    \frac{D h_f^{\circ *}}{D t^*}
    =
    \frac{D }{D t^*}
    \sum_i Y_i^* h_{si}^*(T^*)
    +
    \frac{D }{D t^*}
    \sum_i Y_i^* h_{fi}^{\circ*}
    =
    \\
    \sum_i  
    \frac{D Y_i^*}{D t^*}
    h_{si}^*(T^*)
    +
    \sum_i  
    \frac{D h_{si}^*(T^*)}{D t^*}
    Y_i
    +
    \sum_i  
    \frac{D Y_i^*}{D t^*}
    h_{fi}^{\circ*}
    +
    \cancelto{0}{
    \sum_i  
    \frac{D h_{fi}^{\circ*}}{D t^*}
    Y_i
    }
    =
    \\
    \frac{D }{D t^*}
    \sum_i Y_i^* h_{si}^*(T^*)
    +
    \sum_i  
    \frac{D Y_i^*}{D t^*}
    h_{fi}^{\circ*}
\end{split} 
\end{equation} 

In this case special case $\dfrac{D}{D t^*}=u^*\dfrac{\partial }{\partial x^*}+v^*\dfrac{\partial }{\partial y^*}$ and 
the last term vanish once the formation enthalpy is constant for each specie. 
The term $\dfrac{D Y_i^*}{D t^*}$ is the formation rate of the specie $i$, normally 
write as $\dot{\omega}_i^*$.

%The species equation can multiply by $h_{si}$ and  summed: 
%
%\begin{equation}
%    \rho^*
%    \sum_i  
%    h_{si}^*
%    \frac{D Y_i^*}{D t^*}
%    =
%    \sum_i  
%    \frac{\partial }{\partial y^*}
%    \left(
%        \rho^*D_i^*
%        \frac{\partial Y_i^* }{\partial y^*}
%    \right)
%\end{equation}

So,  with the above sepation the  energy enthalpy equation gets :

\begin{equation}
\frac{D h_{s}^*(T^*)}{D t^*}
=
\mu^*\frac{\partial^2 u^*}{\partial (y^*)^2}+
\frac{\partial}{\partial y^*}
\left(
    k^*\frac{\partial T^*}{\partial y^*}
\right)
-
\frac{\partial}{\partial y^*}
\left(
        \sum_i
        h_{si}
        \rho 
        D_i
        \frac{\partial Y_i^*}{\partial y^*}
\right)
-
\sum_i  
\dot{\omega}_i
h_{fi}^{\circ*}
\end{equation} 

The heat realease of the chemical reaction  
can be writen as: 

\begin{equation} 
    Q_{rxn}
    \dot{\omega}
    =
    \sum_i h_{fi}^{*\circ}
    \dot{\omega}_i
\end{equation} 

where the formation enthalphy  $h_{fi}^{*\circ}$ 
of the specie $i$ is
the enthalphy change when 1 mole of a compound 
is formed from its elements in their standart state. 

For the boundary layer equation the energy equation 
for the reacting misture is  

\begin{equation}
\rho^* u^* \frac{\partial h_s^*}{\partial x^*}+
\rho^* v^* \frac{\partial h_s^*}{\partial y^*}
=
\mu^*\frac{\partial^2 u^*}{\partial (y^*)^2}+
\frac{\partial}{\partial y^*}
\left(
    k^*\frac{\partial T^*}{\partial y^*}
\right)
-
\frac{\partial}{\partial y^*}
\left(
        \sum_i
        h_{si}(T^*)
        \rho^* 
        D_i^*
        \frac{\partial Y_i^*}{\partial y^*}
\right)
-
\sum_i  
\dot{\omega}_i
h_{fi}^{\circ*}
\end{equation} 

% Reaction rate sign convention
\begin{equation} 
\dot{\omega}_i^* > 0 \quad \text{means species } i \text{ is produced}
\end{equation} 

Products of an exothermic reaction usually have lower formation enthalpy

% Standard reaction enthalpy
\begin{equation} 
\Delta H^{\circ*}_{\text{reaction}}
=
\sum_{\text{products}} h_{fi}^{\circ*}
-
\sum_{\text{reactants}} h_{fi}^{\circ*}
< 0
\quad \text{(exothermic reaction)}
\end{equation} 

then the term: 

\begin{equation} 
% Source term in sensible enthalpy equation
\sum_i h_{fi}^{\circ}\dot{\omega}_i < 0
\end{equation} 

if reaction is exothermic. Therefore in the enthalply equation :

\begin{equation} 
- \left( \sum_i h_{fi}^{\circ}\dot{\omega}_i \right) > 0
\end{equation} 

Meaning that positive heat release term
\begin{equation} 
\text{Exothermic reaction} \;\Rightarrow\;
\text{positive contribution to the energy equation}
.
\end{equation} 

Finally, complete the system with the ideal gas law 

\begin{equation} 
p^*=\rho^*R^*T^*
\end{equation} 



